% =====================================================================
\section{Resumen operativo}
\label{sec:resumen}
% =====================================================================

\begin{table}[ht]
\centering
\caption{Teoremas del tema y cuándo usarlos.}
\label{tab:resumen}
\small
\renewcommand{\arraystretch}{1.75}
\begin{tabular}{@{}l l l@{}}
\toprule
\color{cian}Magnitud & \color{cian}Ecuación & \color{cian}Se conserva si\\
\midrule
CM & $\Rc=\frac1M\sum m_{\alpha}\vb{r}_{\alpha}=\frac1M\int\vb{r}\,\dd m$ & ---\\
Momento lineal & $\dot{\vb{P}}=M\ddot{\Rc}=\vb{F}\ext$ & $\vb{F}\ext=\vb{0}$ (o una componente)\\
Momento angular & $\dot{\vb{L}}_{Q}=\vtorque_{Q}\ext$ ($Q$ fijo o CM) & $\vtorque_{Q}\ext=\vb{0}$\\
König I & $\vb{L}_{O}=\Rc\times M\Vc+\vb{L}'$ & ---\\
König II & $T=\frac12MV^{2}+T'$ & ---\\
Energía & $E=T+U\inn+V\ext$ & fuerzas conservativas\\
Masa variable & $m\dot{\vb{v}}=\vb{F}\ext+\dot m\,\vb{u}$ & ---\\
Virial & $2\prom{T}=-\prom{\sum\vb{r}_{\alpha}\cdot\vb{F}_{\alpha}}=n\prom{V}$ & movimiento acotado\\
Pappus & $A=2\pi\bar y L$,\ \ $V=2\pi\bar y A$ & ---\\
\bottomrule
\end{tabular}
\end{table}

\begin{clave}{Estrategia general}
\begin{enumerate}[leftmargin=1.4em]
  \item \emph{Elegir el sistema}: qué partículas entran (en masa variable,
        cuerpo $+$ masa intercambiada; en una explosión, todos los fragmentos).
  \item \emph{Clasificar las fuerzas}: las internas no cambian $\vb{P}$ ni, si son
        centrales, $\vb{L}$; sí pueden cambiar $T$.
  \item \emph{Buscar conservaciones}: $\vb{F}\ext=\vb{0}\Rightarrow\vb{P}$;
        $\vtorque\ext=\vb{0}\Rightarrow\vb{L}$; todo conservativo $\Rightarrow E$.
  \item \emph{Pasar al CM} cuando simplifique: allí $\vb{P}'=\vb{0}$, y por
        König $\vb{L}$ y $T$ se separan en una parte del CM y otra interna.
\end{enumerate}
\end{clave}

\begin{clave}[violeta]{Errores típicos}
\begin{itemize}[leftmargin=1.3em]
  \item Usar $\vb{F}=m\vb{a}$ con masa variable: hay que añadir el empuje
        $\dot m\,\vb{u}$, con $\vb{u}$ \emph{relativa} al cuerpo.
  \item Conservar $T$ en una explosión o un choque: se conserva $\vb{P}$, no $T$.
  \item Olvidar que $\dot{\vb{L}}_{Q}=\vtorque_{Q}\ext$ exige $Q$ fijo o el CM.
  \item En el CM de un continuo, no aprovechar la simetría, o poner la
        densidad sin el factor geométrico ($\dd S=R^{2}\sin\theta\,\dd\theta\,\dd\varphi$).
\end{itemize}
\end{clave}
